\section{A reciprocity obstruction on the complete source curve}
\label{sec:hasse-families}

We now retain the fourth square in the ordered distance construction.
For $p\in\Q$, $p>1$, put $a=p-1$ and $b=p+1$.  The complete source
is the intersection of three quadrics
\begin{equation}\label{eq:hp-full}
 C_p:\quad pV^2=2a(U^2-pZ^2),\qquad
 pW^2=2b(U^2+pZ^2),\qquad pT^2=aU^2+pbZ^2
 \quad\subset\PP^4_{[Z:U:V:W:T]}.
\end{equation}
On $Z\ne0$, the parameter $\rho=U^2/(pZ^2)$ gives the four ordered
squares
\begin{equation}\label{eq:hp-four}
 p\rho,\quad 2(p-1)(\rho-1),\quad
 2(p+1)(\rho+1),\quad (p-1)\rho+p+1.
\end{equation}
The smooth projective model has genus five: its degree-sixteen map to
$\PP^1_\rho$ has five branch values
$0,1,-1,-b/a,\infty$, with inertia of order two at each.
Thus $2g-2=-32+5\cdot8=8$.
Write $\operatorname{Loc}(p)$ for the assertion that the four expressions
in \eqref{eq:hp-four} are simultaneous nonzero squares at a finite
$\rho\notin\{0,1,-1\}$ over each completion of $\Q$.
The real solutions can be chosen in the strict distance region: take
$\rho>1$ sufficiently close to $1$.  Indeed, if
$F=a\rho+b$, $u^2=p\rho$, $v^2=2a(\rho-1)$ and
$w^2=2b(\rho+1)$, the strict margin is
$2F-4|u|-|vw|$, whose limit is $4p-4\sqrt p>0$.

The obstruction below uses every source point on \eqref{eq:hp-full}.
Its formula depends only on the valuations of $p,p-1,p+1$.
We use the additive Hilbert symbol notation $(c,d)_v\in\operatorname{Br}(\Q_v)[2]$,
and write its bit as $2\operatorname{inv}_v(c,d)_v\in\F_2$.
For an odd prime $q$, and $c=q^r c_0,d=q^s d_0$ with units $c_0,d_0$,
\begin{equation}\label{eq:hp-hilbert}
 (-1)^{2\operatorname{inv}_q(c,d)_q}
 =(-1)^{rs(q-1)/2}\leg{c_0}{q}^{s}\leg{d_0}{q}^{r}.
\end{equation}
For odd $u,v\in\Z_2^\times$, the dyadic rule is
\begin{equation}\label{eq:hp-hilbert2}
 2\operatorname{inv}_2(2^r u,2^s v)
 =\frac{(u-1)(v-1)}4+r\frac{v^2-1}8+s\frac{u^2-1}8\pmod2.
\end{equation}

\subsection{The obstruction and its arithmetic domain}

Assume
\begin{equation}\label{eq:hp-domain}
 p(p^2-1)=nk^2,\qquad n>1\text{ odd and squarefree},\qquad k\in\Q^\times.
\end{equation}
For parameters satisfying $\operatorname{Loc}(p)$ define
\begin{align}
 A(p)&=\sum_{\substack{q\text{ odd},\ (2/q)=-1\\v_q(a)>0}}\frac{v_q(a)}2,
 &B(p)&=\sum_{\substack{q\text{ odd},\ (2/q)=-1\\v_q(b)>0}}\frac{v_q(b)}2
 \quad\text{in }\F_2.\label{eq:hp-AB}
\end{align}
If $v_q(p)<0$ is odd, choose $e\in\F_q$ with $e^2=2$ and define
$\sigma_q\in\F_2$ by
\[
 (-1)^{\sigma_q}=\leg{2+e}{q},\qquad
 P(p)=\sum_{\substack{q\text{ odd}\\v_q(p)<0\text{ odd}}}\sigma_q.
\]
This definition is independent of the choice of $e$, since
$(2+e)(2-e)=2$ is a nonzero square.  The existence of $e$ and the
integrality of the half-valuations follow from the next lemma.
Put
\begin{equation}\label{eq:hp-parity}
 \mathcal T(p)=
 \begin{cases}
 B(p)+P(p),&v_2(p)\ne0,\\
 [v_2(a)=5]+B(p)+P(p),&p\equiv1\pmod {16},\\
 [v_2(b)=5]+A(p)+P(p),&p\equiv15\pmod {16}.
 \end{cases}
\end{equation}
Here brackets denote an indicator, and unit congruences are interpreted
in $\Z_2^\times$.

\begin{theorem}\label{thm:hp-uniform}
Suppose \eqref{eq:hp-domain} holds.  If $\operatorname{Loc}(p)$ holds
and $\mathcal T(p)=1$, then $C_p(\Q)=\varnothing$.
\end{theorem}

We prove the theorem by constructing a quaternion algebra on a
coordinate-omission curve and calculating its invariant at every place.
The sum is $\mathcal T(p)/2$.  A rational point would make this sum zero
by Hilbert reciprocity.

Omitting $U$ and $W$, respectively, gives the genus-one curves
\begin{align}
 D_U:&\quad v^2=2t^2-4pz^2,\qquad
 aw^2=2b(t^2-2z^2),\label{eq:hp-DU}\\
 D_W:&\quad V^2=2T^2-4pZ^2,\qquad
 aU^2=pT^2-pbZ^2.\label{eq:hp-DW}
\end{align}
These are projective equations in the four displayed coordinates.
Both omission maps are everywhere defined: if all retained coordinates
vanish, \eqref{eq:hp-full} forces the omitted coordinate to vanish too.

\begin{lemma}\label{lem:hp-norms}
Under $\operatorname{Loc}(p)$, the numbers $p,a,b$ are norms from
$\Q(\sqrt2)$.  At an odd prime where $2$ is nonsquare their valuations
are even.  Under \eqref{eq:hp-domain}, the possible dyadic branches are
\[
 v_2(p)\ge4\text{ even},\quad v_2(p)\le-4\text{ even},\quad
 p\equiv1\pmod {16},\quad p\equiv15\pmod {16}.
\]
In the last two branches, respectively, $[a]_2=[2]_2$ and
$[b]_2=[2]_2$, and the positive valuation is odd and at least five.
\end{lemma}
\begin{proof}
At a completion where $2$ is nonsquare, a point of \eqref{eq:hp-DU}
has $z,w\ne0$.  Its equations give the exact norm witnesses
\[
 \Nm\left(\frac{v+t\sqrt2}{2z(1+\sqrt2)}\right)=p,
 \qquad
 \Nm\left((2+\sqrt2)\frac b w(t+z\sqrt2)\right)=ab.
\]
At an odd inert prime the valuation of a norm is even.  At a pole of
$p$, the valuations of $p,a,b$ agree.  Away from poles, $a,b$ are
integral and cannot both have positive valuation.  The even valuation
of $ab$ therefore makes each positive valuation even.  Thus
$(2,a)_q=(2,b)_q=0$ at every odd place; their real invariants also vanish.
Reciprocity gives vanishing at $2$.  The quadratic norm criterion gives
the asserted global norms.

The norm classes at $2$ are $1,-1,2,-2$.  Since
$v_2(pab)$ is even, a nonzero $v_2(p)$ is even.  Depth $2$ is excluded
by the unit part of $a$: if $p$ has depth $2$ it is $3\pmod8$, and
if $p=u/4$ with $u\equiv1,7\pmod8$, it is $5,3\pmod8$.
For a unit $p$, its norm condition gives $p\equiv1,7\pmod8$;
the norm conditions on $b,a$ sharpen these to $1,15\pmod {16}$.
We spell out the four-square restriction at $2$.
Write $A=v_2(p)$ and $R=v_2(\rho)$.  Since $p\rho$ is square,
$A+R$ is even.  If $A,R$ are both nonzero and have the same sign,
the valuation of $2(p-1)(\rho-1)$ is odd.  If their signs are opposite
and both are even, the same conclusion holds.  If their signs are
opposite and both are odd, the fourth radicand has a unique term of
negative odd valuation: it is $(p-1)\rho$ when $A>0$, and
$p(\rho+1)$ when $A<0$.  Finally $A=R=0$ is impossible: the first
square gives $p\equiv\rho\pmod8$, so one of the middle two radicands
has valuation three.  Exactly one of $A,R$ is therefore zero and the
other is even.

For a unit $p$, the middle two squares force both $v_2(p-1)$ and
$v_2(p+1)$ odd, so $p\equiv1$ or $7\pmod8$.  If $R>0$, the fourth
square excludes $p\equiv1\pmod8$, because its valuation would be one.
Thus $[p]_2=-1$; the product of the middle two squares gives
$[p^2-1]_2=[\rho^2-1]_2=-1$.
If $R<0$, the fourth square excludes $p\equiv7\pmod8$, because
$(p-1)\rho$ would have odd valuation.  Now $[p]_2=1$ and
$[p^2-1]_2=[\rho^2-1]_2=1$.
In both cases $[pab]_2=1$.
For $p\equiv1\pmod {16}$, $[p]_2=1$ and $[b]_2=2$, whence
$[a]_2=2$; for $p\equiv15\pmod {16}$, $[p]_2=-1$ and
$[a]_2=-2$, whence $[b]_2=2$.  The corresponding valuation is odd
and at least five.
\end{proof}

\subsection{A rational anchor and the quaternion class}

Temporarily suppose an actual rational point
$Q_0=[z:v:w:t]\in D_U(\Q)$ is given.  Its $z,w,t$ are nonzero.
For instance, $z=0$ gives $v^2=2t^2$, and $w=0$ gives $t^2=2z^2$;
$t=0$ is excluded over $\R$.  Also $v\ne0$: otherwise
$2p=(t/z)^2$, contrary to the even valuation of $pab$ at $2$.
Define
\begin{equation}\label{eq:hp-anchor}
 z_1=aw-2bz,\qquad u_1=pb(w+2t+2z),\qquad
 t_1=-b(aw+2pt+2az).
\end{equation}
This is a rational point on $aU^2-pT^2+pbZ^2=0$, since
\[
 au_1^2-pt_1^2+pbz_1^2
 =2p^2b(aw^2-2bt^2+4bz^2)=0.
\]
The linear transformation $[z:t:w]\mapsto[z_1:u_1:t_1]$ has determinant
$-8p^2b^2$ and hence never produces the zero tuple.  Choose the sign of
$w$ so that $z_1t-zt_1\ne0$.  Such a choice exists: the sum for the two
signs is $4abz(t+z)$, and $t=-z$ is impossible in \eqref{eq:hp-DU}.
Fix this choice.  On $D_W$ put
\begin{equation}\label{eq:hp-polar}
 L_U=vV-2tT+4pzZ,\qquad
 L_V=au_1U-pt_1T+pbz_1Z,\qquad
 F=\frac{L_UL_V}{V^2},\qquad \alpha=(-a,F).
\end{equation}

Here is the divisor calculation underlying this construction.
The projections from $D_W$ to the two conics in \eqref{eq:hp-DW} have
degree two.  Let $E_U,E_V$ be the degree-two fibres over their respective
rational anchors $[z:v:t]$ and $[z_1:u_1:t_1]$, and let $H_V$ be the
degree-four coordinate section $V=0$.  A tangent line on a smooth
conic meets it twice at its point of tangency.  Consequently
\begin{equation}\label{eq:hp-divisor}
 \operatorname{div}(F)=2(E_U+E_V-H_V).
\end{equation}
All residues of $(-a,F)$ vanish, so it defines a class on the smooth
proper curve $D_W$.

For completeness, the associated elliptic curve and ordered classes are
\begin{equation}\label{eq:hp-elliptic}
 E:\ y^2=x(x-L)(x-R),\qquad
 L=\frac{8p}{a^2},\quad R=-\frac{8pb}{a^3},\qquad
 \chi=(2,pa,2pa),\quad\beta=(-a,-a,1).
\end{equation}
The row $\chi$ describes the actual curve $D_W$.  Indeed, with
$\xi_0=2aU,\xi_L=2V,\xi_R=2T,\zeta=a^2Z$ its equations become
\[
 2\xi_0^2-pa\xi_L^2=L\zeta^2,
 \qquad 2\xi_0^2-2pa\xi_R^2=R\zeta^2.
\]
In particular its degree-four covering map is
$x=8U^2/(a^2Z^2)$, $y=16pUVT/(a^4Z^3)$.
Write $T_0,T_L,T_R$ for the named two-torsion points, and
$\sigma_i$ for the sign change of a coordinate $i$.
The diagonal degree-four covering map is the intrinsic map
$P\mapsto[P-\sigma_Z(P)]$ to the Jacobian: its ordered root functions
are $2U/(aZ)$, $2V/(a^2Z)$ and $2T/(a^2Z)$, with the constants in
$\chi$.  It sends the boundary sections $U=0,V=0,T=0,Z=0$ to
$T_0,T_L,T_R,0$, respectively.
Choose a geometric point $O_C$ above $Z=0$ to compute Abel sums.
Then $\sigma_Z$ acts by negation and the covering map acts by
multiplication by two.  Each other reflection acts as
$P\mapsto t_i-P$, with $2t_i=0$ because it commutes with $\sigma_Z$.
At a fixed point of $\sigma_i$, the covering image is $2P=t_i$;
the boundary images therefore give
$t_U=T_0,t_V=T_L,t_T=T_R$.
Every fibre of the corresponding conic projection has Abel sum
$P+\sigma_i(P)=t_i$, including ramified fibres counted twice.
Thus the Abel sums of $E_U,E_V$ are $T_0,T_L$.
Since $H_V\sim2E_V$, we obtain
$[E_U+E_V-H_V]=T_0-T_L=T_R$.

For the cyclic algebra in \eqref{eq:hp-polar}, the divisor identity
\eqref{eq:hp-divisor} identifies the Picard cocycle as
$\gamma\mapsto\epsilon_{-a}(\gamma)T_R$:
over $\Q(\sqrt{-a})$ the norm of $E_U+E_V-H_V$ is its double,
which is $\operatorname{div}(F)$.
The Weil pairings of $T_R$ with $T_0,T_L,T_R$ are $-1,-1,1$,
respectively, so the ordered two-descent row is $\beta=(-a,-a,1)$.

We shall use only the constant-evaluation consequence of the Picard
cocycle.  Let $k=\Q_v$ and choose any point of $D_W(k)$, which exists
by $\operatorname{Loc}(p)$.  This point splits the degree map on
$\Pic(\overline{D_W})$, giving the Galois module decomposition
$\Pic(\overline{D_W})=E(\overline k)\oplus\Z$.
The coefficient $-a$ becomes a square over $\overline k$, so $\alpha$
lies in the algebraic Brauer group
$\operatorname{Br}_1(D_W)=\ker(\operatorname{Br}(D_W)\to
\operatorname{Br}(\overline{D_W}))$.
The Hochschild--Serre exact sequence
\cite[\S2, (2.1)]{creutz2020} contains
\[
 \operatorname{Br}(k)\longrightarrow\operatorname{Br}_1(D_W)
 \longrightarrow H^1(k,\Pic(\overline{D_W})).
\]
A Kummer class $\beta\in E(k)/2E(k)\subset H^1(k,E[2])$ maps to zero
in $H^1(k,E)$ by the multiplication-by-two exact sequence.  Since
$H^1(k,\Z)=0$, an algebraic Brauer class with this Picard cocycle
comes from $\operatorname{Br}(k)$ and has constant evaluation.
For a torsor the cocycle is computed modulo its torsor class; that
class becomes zero as soon as a local point is chosen.  Thus it
suffices to prove local solubility of the full cover representing
$\beta$.

\subsection{The partner cover at every place}

Put $c=-a$.  The complete partner cover is
\begin{equation}\label{eq:hp-partner}
 cA^2-cB^2=L\zeta^2,\qquad cA^2-C^2=R\zeta^2.
\end{equation}
Two useful projective points are
\begin{align}
 &[1:1:\sqrt c:0] &&\text{if $c$ is square},\label{eq:hp-partner-O}\\
 &\left[\frac{4p}{a^2}\sqrt{\frac b{2p}}:\frac{4p}{a^2}:0:1\right]
 &&\text{if $b/(2p)$ is square}.\label{eq:hp-partner-R}
\end{align}
These follow from $(R-L)/c=(4p/a^2)^2$.

At an odd prime with $v_q(p)>0$, $c$ is square; with $v_q(a)>0$,
$b/(2p)$ is square.  At a pole with $2$ square, the latter criterion
again applies.  At a pole with $2$ nonsquare, write
$p=\pi^{-2h}p_0$, $a=\pi^{-2h}a_0$, $b=\pi^{-2h}b_0$.
Replacing $(A,B,C)$ by $(\pi^{2h}A',\pi^{2h}B',\pi^hC')$ gives a
unit model whose three elliptic roots reduce to $0,8/p_0,-8/p_0$.
It is a smooth genus-one curve over $\F_q$, so the Hasse bound and
Hensel's lemma supply a point.  The same argument applies whenever
$p,a,b$ are units.

It remains to consider $v_q(b)>0$.  If $2$ is square, $c$ is square.
Otherwise write $v_q(b)=2h$ and let $B_0$ be the unit squareclass of
$\pi^{-2h}b$.  The Kummer row at $T_0$ is $(2B_0,2,B_0)$.
If $B_0=1$ this is $\beta$.  If $B_0=2$, choose a residue $X$ with
$X$ nonsquare and $X+2$ a nonzero square.  There are
$(q+\leg{-2}{q})/4>0$ such residues: expanding the two character
indicators reduces the count to
$\sum_X\leg{X(X+2)}q=-1$, which follows from
$(X+1-Y)(X+1+Y)=1$.
Lifting $x=X$ gives an actual point of $E(\Q_q)$ with row $(2,1,2)$,
since $L\equiv-2,R\equiv0$.  Adding $T_0$, whose row is now $(1,2,2)$,
gives $\beta$.  The simultaneous square roots in the Kummer row give
one point of \eqref{eq:hp-partner}.

At $2$, the branches of Lemma~\ref{lem:hp-norms} suffice.
If $v_2(p)\ge4$, $c$ is square.  If $v_2(p)=-2h\le-4$, normalize
$z=1$ in the actual rational anchor.  Its norm condition leaves
$[p]_2=1$ or $-1$.  In the first case the equation
$v^2=2t^2-4p$ forces $v_2(t)=1-h$: a smaller valuation has odd leading
depth and a larger one gives squareclass $-1$.
Then $(2b/a)t^2(1-2/t^2)$ has squareclass $2$, contradicting the
$w$-equation.  Thus $[p]_2=-1$ and $c$ is square.
If $p\equiv1\pmod {16}$, $b/(2p)\in1+8\Z_2$, so
\eqref{eq:hp-partner-R} applies.

Finally suppose $p\equiv15\pmod {16}$.  Normalizing $z=1$, the
first anchor equation forces $v_2(t)\ge2$; the second then implies
that $2b$ is square.  Choose local $r,s,h_d,h_Z$ with
\[
 r^2=2b,\quad s^2=3-p,\quad h_d^2=-a/2,\quad
 h_Z^2=\frac{p^2-4p-1}{2pa}.
\]
Here $3-p=4(1-b/4)$ and $1-b/4\in1+8\Z_2$, so the indicated $s$ exists.
The last two unit radicands are squares because
$-a/2\equiv1\pmod8$ and
$h_Z^2-1=-b^2/(2pa)\in2^8\Z_2$.
The actual point
\[
 [Z:U:V:T]=[h_Z/2:r/4:1:rs/(4h_d)]\in D_W(\Q_2)
\]
has row $\chi=(2,2,1)=\beta$ over $\Q_2$ and hence supplies a point
of \eqref{eq:hp-partner}.  Over $\R$ one may use
$[0:\sqrt{8p/a^3}:\sqrt{8pb/a^3}:1]$.
The partner is therefore everywhere locally soluble, and the evaluation
of $\alpha$ is constant at every place.

\subsection{Evaluation of the quaternion algebra}

We give the local calculations leading to \eqref{eq:hp-parity}.
At a tangent fibre the extension in \eqref{eq:hp-divisor}, rather than
the value of a vanishing rational function, is used.  When the missing
coordinate is locally a square, the actual points
\[
 P_U=[z:u_0:v:t],\quad u_0^2=\frac p a(t^2-bz^2),\qquad
 P_V=[z_1:u_1:v_1:t_1],\quad v_1^2=2t_1^2-4pz_1^2
\]
satisfy
\begin{equation}\label{eq:hp-tangent-eval}
 \alpha(P_U)=(-a,-pL_V(P_U)),\qquad
 \alpha(P_V)=(-a,-2bL_U(P_V)).
\end{equation}
To see this, the polar identities are
\[
 (zV-vZ)^2-2(zT-tZ)^2=-2zZL_U,
 \qquad
 a(z_1U-u_1Z)^2-p(z_1T-t_1Z)^2=-2z_1ZL_V.
\]
Dividing by the square of $tZ-zT$ in the first case gives the limiting
coefficient $-4p/v^2$ for $L_U$.  Dividing by the corresponding square
in the second gives $-p^2b/(2au_1^2)$ for $L_V$; removing squares and
norms from $\Q_v(\sqrt{-a})$ gives the factor $-2b$ in
\eqref{eq:hp-tangent-eval}.  If $u_1=0$, use $U$ as parameter and
$L_V/U^2\to-a/2$.  If $z_1=0$, use the chart $T=t_1$.
The choice $z_1t-zt_1\ne0$ ensures the other polar is nonzero.

For an odd prime the evaluation table is
\begin{equation}\label{eq:hp-odd-table}
\begin{array}{c|c}
 \text{parameter branch}&2\operatorname{inv}_q\alpha\\ \hline
 v_q(p),v_q(a),v_q(b)=0&0\\
 v_q(p)>0&0\\
 v_q(a)>0&2\operatorname{inv}_q(-a,-2)\\
 v_q(b)>0,\ (2/q)=1&0\\
 v_q(b)>0,\ (2/q)=-1&v_q(b)/2\\
 v_q(p)<0\text{ even}&0\\
 v_q(p)<0\text{ odd}&\text{bit of }\leg{-(2+e)}q,\quad e^2=2.
\end{array}
\end{equation}
Here and below equality of entries is in $\F_2$.
We justify the entries, including arbitrary valuation depths.
At a good unit prime, the source and its two conics have smooth proper
models.  Normalize the anchor primitively.  The anchor transformation
has unit determinant, so both tangent sections are primitive.  Their
horizontal divisors are even and neither has a vertical divisor.
The quaternion extends over the model and specializes to
$\operatorname{Br}(\F_q)=0$.  At a positive-$p$ prime, $-a$ is square.

At a positive-$a$ prime where $2$ is square, use the actual point
$[Z:U:V:T]=[1:\sqrt p:0:\pm\sqrt{2p}]$.
Normalize $\min(v_q(t),v_q(z))=0$ and choose the sign so that $L_U$
is a unit.  Direct substitution of \eqref{eq:hp-anchor} gives
$L_V/L_U\equiv-p^2b\pmod q$; all terms involving $aw$ have positive
valuation because $v_q(a)+2v_q(w)\ge0$.
Removing the square pole $V^2$ leaves the class $-b=-2$.
If $2$ is nonsquare, $v_q(a)$ is even.  In the same normalization
$z_1\equiv-4z$, $t_1\equiv-4t$, and the missing square is
$2t_1^2-4pz_1^2\equiv16v^2$.
Choose $v_1\equiv4v$; then $L_U(P_V)\equiv8v^2$ is a unit.
Both entries of the symbol in \eqref{eq:hp-tangent-eval} have unit
squareclass, so its invariant is zero, as is $(-a,-2)_q$.

At a positive-$b$ inert prime put $v_q(b)=2h$ and normalize
$\min(v_q(t),v_q(z))=0$.  The unramified norm $t^2-2z^2$ gives
$v_q(w)=h$ and $v_q(z)=0$.  The anchor transformation has
$v_q(z_1)=h$, $v_q(t_1)\ge2h$.  Hence the missing square at $P_V$
has class $-p=1$.  Choose its sign so that
$L_U(P_V)/z_1\equiv\pm2v-4z$ is nonzero.
Thus $v_q(L_U(P_V))=h$.  The first entry $-a$ has class $2$, and
\eqref{eq:hp-hilbert} gives the bit $2h+h=h$.
If $2$ is square, $-a$ is square instead.

At a pole $v_q(p)=-2h$ with $2$ nonsquare, normalize $z=1$.
The $w$-equation forbids $v_q(t)<0$, and the $v$-equation then
forces $-p$, and hence $-a$, to be square.  If $2$ is square, set
$Z=\pi^hZ'$.  Writing $p=\pi^{-2h}p_0$, $a=\pi^{-2h}a_0$,
$b=\pi^{-2h}b_0$ gives the good model
\[
 V^2-2T^2+4p_0Z'^2=0,\qquad
 a_0U^2-p_0T^2+p_0b_0Z'^2=0.
\]
Normalize $[v:t:z/\pi^h]$ primitively.  If $t$ is nonunit, $-p_0$
is square.  Otherwise the transformed second anchor, multiplied by
$\pi^{4h}$, reduces to
$[0:p_0^2(w+2t):-p_0^2(w+2t)]$, a primitive point.
The original second polar differs from its primitive-model polar by
$\pi^{-6h}$, a square factor.  The good-model argument therefore
again gives zero.

At an odd-depth pole $v_q(p)=-(2h+1)$, $2$ is square.  Normalize $t=1$.
The first anchor equation forces $v_q(z)\ge h+1$ and
$v,w\equiv\pm\sqrt2$.  On $D_W$ choose
$[Z:U:V:T]=[0:u:e:1]$, where $u^2=p/a$, $u\equiv1$ and
$e^2=2$ has sign opposite to $v$.
Then $L_U\equiv-4$; $L_V$ has valuation $-3(2h+1)$ and leading
unit $2p_0^3(w+2)$.  Thus $[F]=[(-a)(w+2)]$ and
$\alpha=(-a,-(w+2))$.  Formula \eqref{eq:hp-hilbert} gives the last
line of \eqref{eq:hp-odd-table}.

We next calculate the dyadic entries.  In either nonunit branch $-a$
is square, as proved above, and the value is zero.
For $p\equiv1\pmod {16}$ put $h=v_2(a)$, $d=(h-1)/2$ and normalize
$z=1$.  The anchor has $t=2t_0$, $v=2v_0$, with $t_0,v_0$ odd and
$v_0^2=2t_0^2-p$; moreover $v_2(w)=(3-h)/2$.
Set
\[
 C=\frac{aw}{2b},\quad g=\frac{t_1}{2z_1},\quad
 z_1=-2b(1-C),\quad v_2(C)=d.
\]
Then $g$ is odd and
\[
 g-t_0=\frac{a(t_0+1/2)+C(t_0+b/2)}{1-C},\qquad
 v_2(g-t_0)\ge d+1.
\]
Choose $e^2=2g^2-p$ close to $-v_0$, and set
$K=v_0e-2t_0g+p$.  Factoring
$e^2-v_0^2=2(g-t_0)(g+t_0)$ gives
$v_2(e+v_0)\ge d+2$, whence $[K]=[-2]_2$.
The actual point $P_V$ with $v_1=2z_1e$ has $L_U=4z_1K$.
Formula \eqref{eq:hp-tangent-eval} becomes
$\alpha_2=(-2,(1-C)K)$.
By \eqref{eq:hp-hilbert2} the $K$ term has bit one, and the
$1-C$ term has bit one exactly when $d=2$.  Therefore
\begin{equation}\label{eq:hp-dyadic-one}
 2\operatorname{inv}_2\alpha=1+[v_2(a)=5].
\end{equation}

For $p\equiv15\pmod {16}$ put $h=v_2(b)$, $d=(h-1)/2$ and again
$z=1$.  Now $t=2t_0$ with $t_0$ even, $v=2v_0$,
$v_0^2=-p+2t_0^2$, and $v_2(w)=(h+1)/2$.
The transformed anchor has
$v_2(z_1)=(h+3)/2$, $v_2(t_1)=h+2$.
Thus $g=t_1/(2z_1)$ has depth $d$, and $e^2=2g^2-p$ is a square.
Choose its sign so that $v_0e\equiv3\pmod4$.
Since $p\equiv-1\pmod {32}$, the number
$K=v_0e-2t_0g+p$ is $14$ or $2\pmod {16}$; its valuation is one
and its odd part is $7$ or $1\pmod8$.
Hence $(2,K)_2=0$ and \eqref{eq:hp-tangent-eval} gives
\begin{equation}\label{eq:hp-dyadic-minus}
 2\operatorname{inv}_2\alpha
 =2\operatorname{inv}_2(2,w/z)+[v_2(b)=5].
\end{equation}
Indeed $z_1=aw(1-2b/(aw))$, the correction has depth $d$, and
$p,a,b$ are norms from $\Q_2(\sqrt2)$.

The invariant at $\R$ is $1/2$.  To verify the sign in the chosen
rational gauge, note that
\[
 (2pt+2az)^2-a^2w^2
 =2\bigl((p^2+1)t^2+4patz+4paz^2\bigr)>0.
\]
The quadratic is positive definite, and $2pt+2az$ has the sign of $t$.
Thus $t_1$ has the opposite sign to $t$.
Cauchy--Schwarz on the two positive conics shows that $L_U$ and $L_V$
have respective signs $-\operatorname{sgn}(tT)$ and
$-\operatorname{sgn}(t_1T)$ away from their zeros.  Consequently
$F<0$.  Since $-a<0$, the real invariant is $1/2$.
The evaluation at a zero or pole is obtained by continuity of the
proper class, so this value holds on every real component.

\subsection{Reciprocity and the exclusion theorem}

The dyadic ratio in \eqref{eq:hp-dyadic-minus} has an exact global
support formula.  At an odd prime where $2$ is nonsquare,
\begin{equation}\label{eq:hp-wsupport}
 v_q(w/z)=
 \begin{cases}
 -v_q(a)/2,&v_q(a)>0,\\
 v_q(b)/2,&v_q(b)>0,\\
 0,&\text{otherwise}.
 \end{cases}
\end{equation}
For a unit or positive-$p$ branch, normalize $\min(v_q(t),v_q(z))=0$.
The unramified norm $t^2-2z^2$ forces $w$ to be a unit, and $z$
cannot be nonunit, since then $v^2$ has leading squareclass $2$.
For $v_q(a)=2h>0$ normalize $w=1$: the same norm has depth $2h$,
and the $v$-equation forces $v_q(z)=h$.
For $v_q(b)=2h>0$ the normalization above gives $v_q(w)=h$, $v_q(z)=0$.
At a pole, normalize $z=1$.  A negative valuation of $t$ would make
$w^2$ have class $2$; with $t$ integral, $t^2-2$ is a unit, so $w$
is a unit.  These arguments prove \eqref{eq:hp-wsupport} at all depths.
Hilbert reciprocity for the literal rational number $w/z$ now gives
\begin{equation}\label{eq:hp-WR}
 2\operatorname{inv}_2(2,w/z)=A(p)+B(p).
\end{equation}

Apply reciprocity also to the constant symbol $(-a,-2)$.
Its real bit is one; its odd contributions are the positive-$a$
entries of \eqref{eq:hp-odd-table} and, at odd-depth poles, the
character $\leg{-1}q$.
At $2$ its bit is one in the $p\equiv1\pmod {16}$ branch and zero
in the other branches.
Adding this zero global sum to the sum for $\alpha$ cancels the real
entry and every positive-$a$ entry.  At odd poles the remaining
character is $\leg{2+e}q$.  Hence
\[
 2\sum_v\operatorname{inv}_v\alpha
 =2\operatorname{inv}_2\alpha+
   2\operatorname{inv}_2(-a,-2)+B(p)+P(p)
 =\mathcal T(p),
\]
by \eqref{eq:hp-dyadic-one}, \eqref{eq:hp-dyadic-minus} and
\eqref{eq:hp-WR}.  This proves the claimed invariant formula.

To finish Theorem~\ref{thm:hp-uniform}, suppose a rational point of
$C_p$ exists.  Its $U$-omission supplies the single rational anchor
used throughout the construction.  A rational point on $C_p$ has
$Z,V,W,T\ne0$ by the arguments following \eqref{eq:hp-DU}, and
$U=0$ would make $pV^2=-2apZ^2<0$.  It therefore gives the
nondegenerate four squares in \eqref{eq:hp-four}.
The construction applies to this anchor and the image point on $D_W$.
Evaluation there is a quaternion algebra over $\Q$, so the sum of its
local invariants is zero.  This contradicts
$\mathcal T(p)=1$.  The theorem follows.

\subsection{A conic family with nonzero obstruction}

Let
\begin{equation}\label{eq:hp-conic}
 r^2+2s^2=8,\qquad p=3-s^2>1.
\end{equation}
Then $r^2=2(p+1)$ and
$[z:v:w:t]=[s:2(p-1):2r:r]$ is an actual rational point of $D_U$.
Every odd inert prime dividing $p-1$ would make
$s^2=2-(p-1)$ a nonsquare unit; thus $A(p)=0$.
Every pole of $p$ has even depth, so $P(p)=0$.
For rational points of this conic, $v_2(s)=1$ and $v_2(r)\ge3$.
Indeed write $r=R/D$, $s=S/D$ with $\gcd(R,S,D)=1$.
If $D$ were even, $R^2+2S^2=8D^2$ modulo $4$ would force $R,S$
even.  Thus $D$ is odd.  Reduction modulo $8$ gives $R,S$ even;
writing $R=2R_1$, $S=2S_1$ gives $R_1^2+2S_1^2=2D^2$.
Here $R_1$ is even and $S_1$ odd.  With $R_1=2R_2$, the identity
$2R_2^2+S_1^2=D^2$ modulo $8$ forces $R_2$ even.
Consequently $p\equiv-1\pmod {32}$ and
$v_2(p+1)=2v_2(r)-1$.
Formula \eqref{eq:hp-parity} reduces to
\begin{equation}\label{eq:hp-thinbit}
 \mathcal T(p)=[v_2(r)=3].
\end{equation}
Thus every locally soluble complete source in this conic family with
$v_2(r)=3$ is a genus-five failure of the Hasse principle.

\begin{theorem}\label{thm:hp-infinite}
For every integer $j\ge0$, set
\[
 N=3+154j,\quad M=N+2,\quad D=N^2+2M^2,\qquad
 r=\frac{8MN}{D},\quad s=\frac{2(N^2-2M^2)}D,\quad p=3-s^2.
\]
The curves $C_p$ form an infinite family of smooth genus-five curves
with points at every completion of $\Q$, including points in the
strict real distance region, and with no rational point.
\end{theorem}
\begin{proof}
The displayed parametrization satisfies \eqref{eq:hp-conic}.
The integers $M,N,D$ are odd, $\gcd(M,N)=1$ and $v_2(r)=3$.
The ratio $m=M/N$ decreases strictly from $5/3$ towards $1$.
On $1<m\le5/3$, the function
$s=2(1-2m^2)/(1+2m^2)$ is negative with derivative
$-16m/(1+2m^2)^2<0$; thus $p=3-s^2$ decreases as a function of $m$.
It follows that the distinct parameters $p$ increase from
$3719/3481>1$.
Their squarefree kernels in $p(p^2-1)$ are odd, since
$[p(p^2-1)]_2=1$.  Such a kernel cannot equal one:
if $p(p^2-1)=y^2$, the legs $(p^2-1)/y,2p/y$ and hypotenuse
$(p^2+1)/y$ give a rational right triangle of area one, contradicting
Lemma~\ref{lem:int:fermat}.
By \eqref{eq:hp-thinbit}, it remains to prove $\operatorname{Loc}(p)$.
We give the local proof directly.

At $2$ use $\rho=b/(3p+1)$.  Its four radicands are
\begin{equation}\label{eq:hp-rho-minus}
 \frac{r^2}{4}\frac{2p}{3p+1},\qquad
 4\frac{-p(p-1)}{3p+1},\qquad
 r^2\frac{2(2p+1)}{3p+1},\qquad
 4p\rho.
\end{equation}
The three unit factors displayed after their square factors lie in
$1+8\Z_2$, because $p\equiv-1\pmod {32}$.
At an odd prime with $v_q(p)>0$ or $v_q(p-1)>0$, the identities
$r^2=2(p+1)$ or $s^2=3-p$, respectively, make $2$ square locally.
Take $\rho=p$.  The first three radicands are then explicit squares,
and the fourth, $p^2+1$, has square residue $1$ or $2$.
If $v_q(p+1)>0$, use \eqref{eq:hp-rho-minus} again: the first and
third radicands divided by $r^2$ have residues $1/4$ and $1$;
the second divided by $4$ has residue $1$.
At a pole of $p$, the identity $s^2=3-p$ makes $-p$ a local square.
Put $h=1/p$ and $\rho=-(p+1)/(p+3)$.  The radicands become
\[
 (-p)\frac{1+h}{1+3h},\quad
 (-4p)\frac{(1-h)(1+2h)}{1+3h},\quad
 4\frac{1+h}{1+3h},\quad 4\frac{1+h}{1+3h}.
\]
Every remaining factor is a principal unit and hence square.
These witnesses treat all nonordinary odd primes and every depth.

At an ordinary prime $q$, the model has five distinct branch points
and smooth genus-five reduction.  If $q\ge101$, Hasse--Weil gives
\[
 \#C_p(\F_q)\ge q+1-10\sqrt q>0.
\]
Any such point gives a nondegenerate $\Q_q$-point.  At a finite branch
point $\rho_0$, write its one vanishing radicand as $u(\rho-\rho_0)$
with $u$ a unit, and take $\rho=\rho_0+q^2/u$.
The other three radicands remain square units.  At infinity the three
leading ratios to $p\rho$ are squares; taking
$\rho=1/(pq^2)$ gives four nonzero roots.  Thus normalized branch
points are sufficient for lifting.

Here is a short finite-field proof for the remaining primes.  It also
records the explicit arithmetic used for the small-prime part.
Let $\chi$ be the quadratic character of $\F_q$, extended by $\chi(0)=0$,
and put
\[
 H=\chi(-1),\quad T=\chi(2),\quad A=\chi(p),\quad B=\chi(p-1),
 \quad c=-\frac{p+1}{p-1}.
\]
The conic identity gives $\chi(p+1)=T$.
The five elliptic character quotients are
\begin{align*}
 E_0 &:y^2=p(p^2-1)x(x-1)(x+1),\\
 E_1 &:y^2=2p\,x(x-1)(x-c),\\
 E_2 &:y^2=2p(p^2-1)x(x+1)(x-c),\\
 E_3 &:y^2=(p+1)(x-1)(x+1)(x-c),\\
 E_4 &:y^2=p(p+1)c(x-1)(x+1)(x-c^{-1}).
\end{align*}
If $t_i=q+1-\#E_i(\F_q)$, expansion by the characters of the four
square roots gives
\begin{equation}\label{eq:hp-five-traces}
 \#C_p(\F_q)=q+1-\sum_{i=0}^4t_i.
\end{equation}
Indeed the products of three radicands give $E_0,E_1,E_2,E_3$;
the product of all four, after $x=1/\rho$, gives $E_4$.
The remaining characters have genus zero and their finite sums cancel
with their contributions at infinity.

Each $E_i$ has rational two-torsion.  On
$y^2=d\prod(x-e_i)$, the point $(e_i,0)$ is twice a rational point
if and only if $d(e_i-e_j)$ and $d(e_i-e_k)$ are both squares.
This follows by writing the tangent at a putative half-point through
$(e_i,0)$ and factoring the resulting quadratic; conversely those two
square roots give the two tangent slopes and a half-point.
Consequently the group order is $4\pmod8$ when there is no rational
point of order four, is divisible by eight when such a point exists,
and is divisible by sixteen when all four-torsion is rational.

Let $U_q,V_q,W_q$ be the largest integers at most $2\sqrt q$
congruent to $q+1\pmod8$, $q-3\pmod8$, and $q+1\pmod {16}$,
respectively.  Set
$\tau_q=-\sum_{x\in\F_q}\chi(x^3-x)$.
Then $t_0=TAB\tau_q$; for $q\equiv3\pmod4$, $\tau_q=0$ by
pairing $x$ and $-x$.
Applying the preceding halving criterion to the displayed roots gives
the following upper bounds on the sum in \eqref{eq:hp-five-traces}.
Only the sign patterns without an immediate rational branch point
need be listed:
\begin{center}
\begin{tabular}{ccl}
\toprule
$q\bmod8$&sign pattern&upper bound for $\sum t_i$\\
\midrule
$3$&all $A,B$&$3U_q+V_q$\\
$7$&$A=-1,B=1$&$2U_q+2V_q$\\
$1$&$A=B=-1$&$\tau_q+2U_q+2V_q$\\
$5$&$A=-1,B=1$&$\tau_q+4U_q$\\
$5$&$A=B$&$-\tau_q+2U_q+V_q+W_q$\\
\bottomrule
\end{tabular}
\end{center}
For example, in the first row exactly three of $E_1,\ldots,E_4$ have
order divisible by eight, and the remaining order is $4\pmod8$.
In the second and third rows two have each type; in the fourth all
four have order divisible by eight; in the last, one has full rational
four-torsion, two have an order divisible by eight, and one has order
$4\pmod8$.  These statements follow by inserting $A,B,H,T$ into
$d(e_i-e_j)$ and $d(e_i-e_k)$.
The omitted branch patterns are also explicit: when $T=1$, $\rho=1$
works for $A=1$, and $\rho=0$ works for $HB=1$; when $H=1,T=-1$,
$\rho=c$ works for $A=1,B=-1$.

The trace values needed below are
\[
\begin{array}{c|rrrrrrrrrrr}
q&5&13&17&29&37&41&53&61&73&89&97\\\hline
\tau_q&-2&6&2&-10&-2&10&14&-10&-6&10&18.
\end{array}
\]
They can be evaluated from the displayed finite sum; equivalently,
for $q=u^2+v^2$, $u$ odd, $v$ even and $u+v\equiv1\pmod4$, the
usual Gaussian evaluation of this sum is $\tau_q=2u$.
Substitution in the bounds proves positivity except for the following
cases, where the two conic conditions
$\chi(2(p+1))=1$ and $\chi(3-p)\in\{0,1\}$ are retained.
The congruence $N\equiv3$, $M\equiv5\pmod {7\cdot11}$ gives
$p\equiv1\pmod7$ and $p\equiv9\pmod {11}$.
At $17$ the remaining residues are $7,12$, and at $41$ they are
$7,35$; taking $\rho=p$ works because $p^2+1$ has respective square
residues $16,9$ and $9,37$.
At $13$ the sign case $A=-1,B=1$ would require $p=5$, but
$3-p=11$ is nonsquare.
At $53$ the remaining case is $A=-1,B=1$, and at $61$ the case
$A=B=-1$ is also settled by $\rho=-2$, whose radicands are
$-2p,-6(p-1),-r^2,s^2$; use $(3/53)=-1$ and $(3/61)=1$.
The residual cases have the following explicit nonzero square roots:
\begin{center}\small
\begin{tabular}{rrrr}
\toprule
$q$&$p\bmod q$&$\rho\bmod q$&roots in the order \eqref{eq:hp-four}\\
\midrule
11&9&4&$5,2,1,3$\\
19&2&3&$5,2,9,5$\\
19&11&5&$6,2,7,9$\\
19&13&10&$4,8,2,1$\\
19&17&4&$7,1,3,5$\\
23&17&5&$4,6,3,11$\\
29&7&22&$3,7,7,13$\\
29&25&4&$10,12,12,8$\\
43&2&3&$7,2,14,7$\\
43&6&10&$19,2,5,10$\\
43&11&9&$20,17,5,4$\\
43&21&14&$6,2,12,1$\\
43&25&35&$12,16,18,7$\\
43&29&5&$4,3,4,16$\\
43&31&10&$3,14,4,17$\\
43&32&3&$15,9,7,13$\\
43&33&12&$3,4,14,17$\\
43&36&23&$21,11,20,5$\\
61&16&3&$29,11,21,1$\\
61&42&12&$4,29,9,13$\\
61&49&12&$10,18,18,4$\\
\bottomrule
\end{tabular}
\end{center}
Every row is checked by squaring its four entries in \eqref{eq:hp-four};
all roots are nonzero and therefore lift by Hensel's lemma.
This completes the proof of finite-place solubility.
At the real place the strict open interval described after
\eqref{eq:hp-four} supplies the required points.
Now \eqref{eq:hp-thinbit} and Theorem~\ref{thm:hp-uniform} give
$C_p(\Q)=\varnothing$ for every parameter in the family.

There are also infinitely many geometric isomorphism classes among these
curves. A fixed smooth curve of genus five has finite automorphism group:
its automorphism scheme is of finite type, for example through the
tricanonical embedding, and its tangent space is $H^0(C,T_C)=0$ because
$\deg T_C=-8$. It therefore has finitely many subgroups isomorphic to
$(\Z/2)^4$. Each such subgroup gives one quotient with its unordered
branch configuration. Here that configuration is
$\{0,1,-1,c,\infty\}$, with $c=-(p+1)/(p-1)$. For a fixed configuration,
there are finitely many ways to choose and normalize three ordered branch
points, and hence finitely many values of $c$. Since $p=(c-1)/(c+1)$,
a fixed geometric curve occurs for only finitely many of our distinct
parameters.
\end{proof}
